%%% Font, caratteri speciali e lingua %%% 
\usepackage[T1]{fontenc}
\usepackage[utf8]{inputenc}
\usepackage{booktabs}
\usepackage{longtable}
\usepackage[french]{babel}

\usepackage{longtable,array}

%%% Per ringraziamenti a forma di cuore
\usepackage{shapepar}
\usepackage{lipsum} % For placeholder text

%%% simbolo mail
\usepackage{fontawesome}

%\raggedbottom %annulla riempimento della pagina in verticale con spazi aggiuntivi tra capoversi o ambienti

%%% Stile della pagina %%%
\pagestyle{headings} %Ved. p. 35
\usepackage{emptypage}
%%%%%%%%%%%%%%%%

\usepackage{graphicx}
\usepackage[paperheight=9in, paperwidth=6in, twoside]{geometry}
%\usepackage[a4paper]{geometry}
\usepackage{afterpage}
\usepackage[svgnames, dvipsnames]{xcolor}

\usepackage{pagecolor}
%\usepackage{gfsdidot}
\usepackage{charter}
\usepackage{qrcode}
\usepackage{etoolbox}
\renewenvironment{quote}{%
   \list{}{%
     \leftmargin0.1cm   % this is the adjusting screw
     \rightmargin\leftmargin
   }
   \item\relax
}
{\endlist}


\DeclareRobustCommand{\augiefamily}{%
  \fontfamily{augie}\fontseries{m}\fontshape{n}\selectfont}
\DeclareTextFontCommand{\textaugie}{\augiefamily}

\usepackage[object=vectorian]{pgfornament}

\usepackage{changepage}

\usepackage{titling}
\patchcmd{\titlepage}{\setcounter{page}\@ne}{}{\message {Successfully patched titlepage.}}{\message {Failed to patch titlepage.}}


\usepackage[skins]{tcolorbox}

%%%%%%%%%%%%%%%%%%%%%%%%%%
% Dedication's package
\usepackage{emerald}
\pdfmapfile{=emerald.map}
%%%%%%%%%%%%%%%%%%%%%%%%%%

\newtcbox{\whiteshadowbox}[1][]{%
  enhanced, 
  center upper,
  fontupper=\large\bfseries,
  drop fuzzy shadow southeast,
  boxrule=0pt,
  sharp corners,
  colframe=blue,
  colback=white!10,
  before={\begin{center}},
    after={\end{center}},
  #1%
}
%%%%%%%%%%%%%%%%%%%%%%%%%
\usepackage{tocloft}


%%% PACKAGES
\usepackage{dramatist}

\usepackage{titling}% Allows for multiple \titles within one document

\let\oldmaketitle\maketitle
\renewcommand{\maketitle}{%
\newpage
\phantomsection
  \addcontentsline{toc}{chapter}{\thetitle}
\oldmaketitle

\markboth{\MakeUppercase{\thetitle}}{\MakeUppercase{\thetitle}}
}

\renewcommand{\casttitlename}{\includegraphics[scale=.14]{emoji/personadzo.png}\hspace{.5mm} Personadzo \includegraphics[scale=.15]{emoji/personadzo.png}}
\renewcommand{\actname}{}
\renewcommand{\scenename}{\textit{Scène}}
\renewcommand{\printscenenum}{%
\scenenumfont \thescene}

\renewcommand{\StageDir}[1]{%
    \begin{stagedir}
    #1
    \end{stagedir}
    %\vspace*{.5cm}
}

\setlength{\afteractskip}{10pt}
\setlength{\beforesceneskip}{10pt}
\setlength{\aftersceneskip}{10pt}
\setlength{\aftercasttitleskip}{20pt}


%%%%%%%%%%%%%%%%%%%%%%%%
% FOTO COPERTINA
%%%%%%%%%%%%%%%%%%%%%%%%
\newcommand{\fotocopertina}[4]{
\cleardoublepage 
%\newpage
\vspace*{\fill}
\begin{center}
\textbf{\LARGE Le Digourdì di #4}
\end{center}

\begin{figure}[h]
\centering
\whiteshadowbox{
\includegraphics[height=4cm,keepaspectratio]{#1}
}
\end{figure}
\noindent
\textbf{Drette}: \textit{#2.}
\newline
\newline
\textbf{Achouatoù}: \textit{#3.}
}
%%%%%%%%%%%%%%%%%%%%%%%%

%%%%%%%%%%%%%%%%%%%%%%%%
% LINK PIESE
%%%%%%%%%%%%%%%%%%%%%%%%
\newcommand{\LinkPiese}[3]{
\vspace*{\fill}
\begin{figure}[h]
\centering
\foreach \x [count=\i] in {#1} {
 \foreach \y [count=\j] in {#2} {
  \ifnum\i=\j
      \begin{subfigure}{#3\textwidth}
  \centering
    \video\hspace*{0.5mm} \textsc{\x}\hspace*{0.5mm} \video\\
 \vspace*{2mm}
\qrcode[hyperlink, height=0.75in]{\y}
  \end{subfigure}%
   \fi
 }
}
\end{figure}
\vspace*{\fill}
}
%%%%%%%%%%%%%%%%%%%%%%%%

%%%%%%%%%%%%%%%%%%%%%%%%
% Souvenir
%%%%%%%%%%%%%%%%%%%%%%%%
\newcommand{\souvenir}[2]{
\section*{Lo souvenir de l'atteur}
\og #1\fg{}
\newline
\newline
\hspace*{\fill} \textit{#2}
}
%%%%%%%%%%%%%%%%%%%%%%%%

%%%%%%%%%%%%%%%%%%%%%%%%
% Queriaouzitoù
%%%%%%%%%%%%%%%%%%%%%%%%
\newcommand{\queriaouzitou}[1]{
\subsection*{Queriaouzitoù}
#1
}
%%%%%%%%%%%%%%%%%%%%%%%%

%%%%%%%%%%%%%%%%%%%%%%%%
% Scénographie
%%%%%%%%%%%%%%%%%%%%%%%%
\newcommand{\Scenographie}{
\newpage
\begin{center}
\Large
\textsc{\scenographie\ \textit{Scénographie} \scenographie}
\end{center}
}
%%%%%%%%%%%%%%%%%%%%%%%%

%%%%%%%%%%%%%%%%%%%%%%%%
% CAST
%%%%%%%%%%%%%%%%%%%%%%%%
\newcommand{\name}[1]{eunterprétoù pe \textsc{#1}.}
\newcommand{\nameF}[1]{eunterprétaye pe \textsc{#1}.}
%%%%%%%%%%%%%%%%%%%%%%%%

%%%%%%%%%%%%%%%%%%%%%%%%
% Dérì le Rido
%%%%%%%%%%%%%%%%%%%%%%%%
\newcommand{\DeriLeRido}{
\newpage
%\cleardoublepage 
\thispagestyle{empty}
\begin{center}
\textbf{\LARGE Derì le rid\'o}
\end{center}
\vspace{.2cm}
}

\usepackage{pgffor}
\newcommand{\RoleNoms}[2]{
\begin{center}
 \textbf{#1}\\\vspace{.05cm}
  \foreach \x in {#2} {
    \x\\
  	}
  	\vspace{.01cm}
\end{center}
}
%%%%%%%%%%%%%%%%%%%%%%%%  

%%%%%%%%%%%%%%%%%%%%%%%%
% Foua vouése
%%%%%%%%%%%%%%%%%%%%%%%%
\newcommand{\Fv}[1]{
Vouése derì le rid\'o \vouese\ : \og #1 \fg . 
}
%%%%%%%%%%%%%%%%%%%%%%%%

%%%%%%%%%%%%%%%%%%%%%%%%
% VIDEO
%%%%%%%%%%%%%%%%%%%%%%%%
\usepackage{float}
\newcommand{\listvideos}{Liste di \textit{vidéo}}
\newlistof{videos}{vds}{\listvideos}

\newcommand{\StartVideo}[3]{
\refstepcounter{videos}

\begin{figure}[H]
%\vspace*{-5pt}
\StageDir{\hspace{2.5em}Teuppe \lemieBa .}
\centering
\foreach \x [count=\i] in {#1} {
 \foreach \y [count=\j] in {#2} {
  \ifnum\i=\j
      \begin{subfigure}{.75\textwidth}
  \centering
    \video\hspace*{0.5mm} \textsc{\small\y}\hspace*{0.5mm} \video\\\vspace*{2mm}
    \qrcode[hyperlink, height=0.5in]{\x}
  \end{subfigure}%
  \ifnum\i=1
   \vspace{2.5mm}
   \fi
  \addcontentsline{vds}{section}{\y}
   \fi
 }
}
\ifstrempty{#3}%
    {\StageDir{Lemie \lemieSi .}
    {}
%\vspace*{-5pt}
}%
\end{figure}%
#3
\ifstrempty{#3}
    {}
    {\vspace*{-2em} \StageDir{Lemie \lemieSi .}
    }
}
%%%%%%%%%%%%%%%%%%%%%%%%
\usepackage{subcaption}
%%%%%%%%%%%%%%%%%%%%%%%%
% MEZEUCCA
%%%%%%%%%%%%%%%%%%%%%%%%
\usepackage{xparse}
\usepackage{xstring}

\newcommand{\listsounds}{Liste di mezeucque}
\newlistof{sounds}{snd}{\listsounds}

\NewDocumentCommand{\sound}{m m O{true}}{%
  \refstepcounter{sounds}

  \begin{figure}[h!]
    \centering
    \soundcloud \hspace*{0.5mm} \href{#1}{\color{RawSienna}{\textsc{#2}}} \hspace*{0.5mm} \soundcloud%
    \vspace{-5pt}
    % \vspace*{2mm}
    % \qrcode[hyperlink, height=0.5in]{#1}
  \end{figure}

  \IfStrEq{#3}{true}{%
    \addcontentsline{snd}{section}{#2}
  }{}
}
%%%%%%%%%%%%%%%%%%%%%%%%

%%%%%%%%%%%%%%%%%%%%%%%%
% Effet
%%%%%%%%%%%%%%%%%%%%%%%%
\newcommand{\listeffets}{Liste di-z-éffé sonore}
\newlistof{effets}{eff}{\listeffets}

\NewDocumentCommand{\effet}{m m O{true}}{%
  \refstepcounter{effets}

  \begin{figure}[h]
    \centering
    \soundcloud \hspace*{0.5mm} \href{#1}{\color{RawSienna}{\textsc{#2}}}\hspace*{0.5mm} \soundcloud%
  \end{figure}%
  \vspace{-0pt}
  \IfStrEq{#3}{true}{%
    \addcontentsline{eff}{section}{#2}%
  }{}
}
%%%%%%%%%%%%%%%%%%%%%%%%

%%%%%%%%%%%%%%%%%%%%%%%%
% Foto
%%%%%%%%%%%%%%%%%%%%%%%%
\newcommand{\listfotos}{Liste di foto}
\newlistof{fotos}{fot}{\listfotos}

\newcommand{\foto}[2]{
\refstepcounter{fotos}

\begin{figure}[h]
\centering
\instagram \hspace*{0.5mm} \textsc{#2}\hspace*{0.5mm} \instagram\\
 \vspace*{1mm}
\qrcode[hyperlink, height=0.5in]{#1}
\end{figure}
\addcontentsline{fot}{section}{#2}
}
%%%%%%%%%%%%%%%%%%%%%%%%

\newcommand{\ridoiver}{
%\vspace*{-.55cm}
\begin{center}
\raisebox{-.2mm}{\includegraphics[scale=.1]{emoji/ridoIver.png}} {\large\textsc{Se ivre lo rid\'o}} \raisebox{-.2mm}{\includegraphics[scale=.1]{emoji/ridoIver.png}}
\end{center}
}

\newcommand{\ridocliou}{\vspace*{15pt}
\begin{center}
\raisebox{.1mm}{\includegraphics[scale=.07]{emoji/tsilateila.png}} {\large \textsc{tsi la tèila}}  \raisebox{.1mm}{\includegraphics[scale=.07]{emoji/tsilateila.png}}
\end{center}
\vspace*{.25cm}
}

\addto\captionsfrench{\renewcommand{\contentsname}%
    {Pièse}%
}
\newcommand{\moublotitlename}{Moublo}
\newcommand{\moublo}{\newpage \centering\textsc{\moublotitlename}}

%%%%%%%%%%%%%%%%
%%% File sty personali  %%%

%%% Emoticons %%%
\usepackage{emoticons}

%%%%%%%%%%%%%%%%

%%% Collegamenti ipertestuali %%%
\usepackage{hyperref}
\hypersetup{
    colorlinks=true,
    linktoc=all,
    linkcolor=Brown,
    urlcolor = Sepia
}



